\documentclass[11pt]{article}

\usepackage[margin=1in]{geometry}
\usepackage{amsmath,amssymb,amsthm}
\usepackage{graphicx}
\usepackage[hidelinks]{hyperref}

\newtheorem{theorem}{Theorem}
\newtheorem{lemma}{Lemma}
\newtheorem{remark}{Remark}

\newcommand{\norm}[1]{\left\lVert #1 \right\rVert}
\newcommand{\rhoBar}{\overline{\rho}}
\newcommand{\deltaBar}{\overline{\delta}}
\newcommand{\wstar}{w^*}

\title{A Sharpened Asplund Averaging Exponent for Simultaneous AUS/AUC Renormings}
\author{Automated Math Discovery Smoke Test}
\date{June 14, 2026}

\begin{document}
\maketitle

\section*{Source and Open Question}

This packet concerns Remark 4.4 from:

\begin{quote}
Stephen J. Dilworth, Denka Kutzarova, Gilles Lancien, Lovasoa N.
Randrianarivony, \emph{Equivalent norms with the property $(\beta)$ of
Rolewicz}, arXiv:1506.07978.
\end{quote}

The source paper proves that if a reflexive Banach space admits separately an
equivalent AUS norm of power type $p$ and an equivalent AUC norm of power type
$q$, then it admits a single equivalent norm which is simultaneously AUS of
power type $p$ and AUC of power type $4q$. Immediately after this lemma, the
authors ask whether both original power types can be preserved.

The source location is Remark 4.4 on page 10.

\begin{center}
\includegraphics[width=\textwidth]{figures/open_problem_crop.png}
\end{center}

\section*{Proof Intuition}

Lemma 4.3 in the source paper uses dual norms
\[
  \norm{x^*}_{n,*}=N^*(x^*)+n^{-1}M^*(x^*)
\]
and averages their predual norms with weights $n^{-3}$. In the AUC part of
the proof, the paper records a strict $t^q$ gain at one index
$n_0\simeq t^{-q}$, producing the displayed loss $t^{4q}$.

The same argument actually gives a strict $t^q$ gain for every tail index
$n\ge n_0$, not just for the single index $n_0$. If the averaging weights are
changed from $n^{-3}$ to $n^{-\theta}$ with $\theta>2$, the AUS estimate is
unchanged because $\sum n^{1-\theta}<\infty$. The AUC estimate gains the whole
tail
\[
  t^q\sum_{n\ge n_0} n^{-\theta}
  \simeq t^q n_0^{1-\theta}
  \simeq t^{q\theta}.
\]
Choosing $2<\theta<r/q$ gives AUC power type $r$ for every $r>2q$.

This still does not answer the exact preservation question: the method cannot
reach $q$, and even this weighted version stops strictly above $2q$.

\section*{Candidate Partial Result}

\begin{theorem}
Let $X$ be a reflexive Banach space. Suppose that $X$ admits an equivalent
AUS norm of power type $p$ and an equivalent AUC norm of power type $q$, where
$1<p,q<\infty$. Then, for every $r>2q$, $X$ admits an equivalent norm which is
simultaneously AUS of power type $p$ and AUC of power type $r$.
\end{theorem}

\begin{proof}
We modify the proof of Lemma 4.3 in Dilworth--Kutzarova--Lancien--
Randrianarivony. Let $N$ be an equivalent norm with
\[
  \deltaBar_N(t)\ge a t^q\qquad (0<t\le 1)
\]
and let $M$ be an equivalent norm with
\[
  \rhoBar_M(t)\le b t^p\qquad (0<t\le 1).
\]
After replacing constants, assume
\[
  N(x)\le M(x)\le C N(x)\qquad (x\in X).
\]

For $n\ge 1$ define an equivalent dual norm on $X^*$ by
\[
  \norm{x^*}_{n,*}=N^*(x^*)+n^{-1}M^*(x^*),
\]
and let $\norm{\cdot}_n$ denote its predual norm on $X$. As in Lemma 4.3 of
the source paper, there is a constant $K$ such that
\[
  \rhoBar_{\norm{\cdot}_n}(t)\le K n t^p
  \qquad (0<t\le 1,\ n\ge 1).                 \tag{1}
\]
Also
\[
  (1+n^{-1})^{-1}N(x)\le \norm{x}_n\le N(x)
  \qquad (x\in X).                            \tag{2}
\]

Fix $r>2q$, and choose $\theta$ with
\[
  2<\theta<r/q.
\]
Define
\[
  |x|=\sum_{n=1}^{\infty} n^{-\theta}\norm{x}_n .
\]
By (2), this is an equivalent norm on $X$.

\subsection*{AUS estimate}

Let $|x|=1$ and $0<t\le 1$. Since $|\cdot|$, $N$, and the norms
$\norm{\cdot}_n$ are uniformly equivalent, there are constants $A,B>0$ such
that
\[
  A|z|\le \norm{z}_n\le B|z|,
  \qquad
  A|z|\le N(z)\le B|z|
\]
for all $z\in X$ and $n\ge 1$.

Choose $n_1$ so large that
\[
  2B\sum_{n>n_1} n^{-\theta}\le t^p .
\]
For each $1\le n\le n_1$, (1) gives a finite-codimensional subspace on which
\[
  \norm{x+y}_n\le \norm{x}_n+K' n t^p
\]
whenever $|y|\le t$, where $K'$ is independent of $n,x,t$. Intersecting these
finitely many subspaces, and then summing, yields
\[
\begin{aligned}
  |x+y|
  &\le \sum_{n\le n_1}n^{-\theta}\norm{x+y}_n
       +\sum_{n>n_1}n^{-\theta}\norm{x+y}_n\\
  &\le \sum_{n\le n_1}n^{-\theta}\norm{x}_n
       +K't^p\sum_{n=1}^{\infty}n^{1-\theta}
       +\sum_{n>n_1}n^{-\theta}\norm{x}_n+t^p\\
  &\le 1+C_1t^p .
\end{aligned}
\]
Here $\sum n^{1-\theta}<\infty$ because $\theta>2$. Thus $|\cdot|$ is AUS of
power type $p$.

\subsection*{AUC estimate}

Let $|x|=1$ and $0<t\le 1$. The AUC assumption for $N$ gives a constant
$\eta>0$ and a finite-codimensional subspace $E_0$ such that, for all
$y\in E_0$ with $|y|\ge t$ and $|y|\le 1$,
\[
  N(x+y)\ge N(x)(1+\eta t^q).                 \tag{3}
\]
Choose a large constant $\alpha$ and an integer $n_0$ satisfying
\[
  n_0-1\le \alpha t^{-q}\le n_0.
\]
Taking $\alpha$ large enough, (2) and (3) imply that, for every $n\ge n_0$,
\[
\begin{aligned}
  \norm{x+y}_n
  &\ge (1+n^{-1})^{-1}N(x+y)\\
  &\ge (1-n^{-1})N(x)(1+\eta t^q)\\
  &\ge \norm{x}_n+d\,t^q
\end{aligned}                                      \tag{4}
\]
with a constant $d>0$ independent of $x,y,t,n$. In the last step we used
$\norm{x}_n\le N(x)$ and the lower bound $N(x)\ge A|x|=A$.

For the finitely many indices $n<n_0$, choose norming functionals for $x$ in
the norms $\norm{\cdot}_n$ and intersect their kernels with $E_0$. The resulting
finite-codimensional subspace $E$ still satisfies (3), and for $y\in E$,
\[
  \norm{x+y}_n\ge \norm{x}_n \qquad (n<n_0).       \tag{5}
\]
Combining (4) and (5), for all $y\in E$ with $|y|\ge t$ and $|y|\le 1$,
\[
\begin{aligned}
  |x+y|
  &\ge |x| + d t^q \sum_{n\ge n_0}n^{-\theta}\\
  &\ge 1+d_1 t^q n_0^{1-\theta}\\
  &\ge 1+d_2 t^{q\theta}.
\end{aligned}
\]
Thus $|\cdot|$ is AUC of power type $q\theta$.

Since $q\theta<r$ and $0<t\le 1$, the estimate
$\deltaBar_{|\cdot|}(t)\ge d_2t^{q\theta}$ implies
$\deltaBar_{|\cdot|}(t)\ge d_2t^r$ after decreasing constants if necessary.
Therefore $|\cdot|$ is AUC of power type $r$, while the first part showed it
is AUS of power type $p$.
\end{proof}

\section*{Verification Notes}

The proof is not computational. It is a direct modification of the source
paper's Asplund averaging construction. The only new point is the tail
summation in the AUC estimate. The proof deliberately keeps the same auxiliary
dual norms $N^*+n^{-1}M^*$ used in the paper, so that the AUS estimate
quoted as (1) is exactly the estimate already used in Lemma 4.3.

The result is a partial answer only. It improves the source-paper exponent
$4q$ to every exponent strictly larger than $2q$, but the question from Remark
4.4 asks for the exact exponent $q$.

\section*{Novelty and Search Bounds}

Before promoting this packet, the run indexes
\[
\texttt{registry\_index.tsv},\quad
\texttt{solutions/index.tsv},\quad
\texttt{attempts/index.tsv},\quad
\texttt{proof\_gaps/index.tsv}
\]
were searched for arXiv:1506.07978, Rolewicz, AUS/AUC power types, and the
phrase ``preserve both power types.'' The local parsed arXiv source corpus was
also searched for the exact phrase and close AUS/AUC power-type formulations.
No duplicate packet or local later answer was found.

A bounded web search for close phrases such as ``preserve both power types,''
``both power types together,'' and ``simultaneously AUS AUC power type'' found
no exact later resolution in that search window. This is not a full literature
survey, so the novelty confidence should be treated as bounded rather than
definitive.

\section*{Human Review Recommendation}

Recommended verdict after review, if the estimates check out: accept as a
partial result improving Lemma 4.3's quantitative exponent loss. The remaining
open question is whether the AUC exponent can be reduced from every $r>2q$ all
the way to $q$.

\begin{thebibliography}{9}
\bibitem{DKLR2015}
S. J. Dilworth, D. Kutzarova, G. Lancien, and L. N. Randrianarivony,
\emph{Equivalent norms with the property $(\beta)$ of Rolewicz},
arXiv:1506.07978, 2015.
\end{thebibliography}

\end{document}
